% Transcription — fonds Alexandre Grothendieck, shelfmark 10, batch 5 (pages 81-100).
% Established from the facsimile archives/batches/10/batch-05.pdf (a mirror of
% https://grothendieck.umontpellier.fr/10.pdf), per the skill
% transcribe-grothendieck.
%
% Pass: Opus 5 (claude-opus-5), 2026-09-12 — first pass, unchecked against the
% pages by a human.
%
% The batch breaks in three. Pages 81-83 finish the typescript of batch 4 and
% stop, mid-run, at a « Cor. » with nothing after it. Pages 85-86 are one leaf
% of an unrelated typescript on classifying topoi, corrected throughout in his
% hand, with a census of the simple types on its back in his hand alone. Pages
% 88-98 are a run of his own, opened by a cover sheet and carried to the end of
% the batch, on filtrations of a fibre functor.
%
% Pages 81-83, the end of « Catégories tannakiennes définies par des
% cristaux » (n° 9 and n° 10, « Cas k fini »):
%
%   81   the slope-[0,1] isocrystals as the Dieudonné modules of p-divisible
%        formal groups, the category being opposite to theirs mod isogeny, and
%        amplitude in the closed [0,1] as the Barsotti–Tate condition, the unit
%        isocrystal answering the ind-étale Q_p/Z_p; then for k perfect the
%        equivalence between effectivity plus a V_i with FV_i = V_iF = p^i·id
%        and amplitude in [0,i]; and the canonical fibre functor over
%        K_a = K(F_{q=p^a}), whose extra structure is the linear automorphism
%        F_E^a, with the open subgroup G_a ⊂ G_{K_a} of index a cut out by the
%        a-th roots of unity in the character group Q*.
%   82   Proposition 10.1: FCriso(F_q)^tann ⊗ K_a ≃ FCriso(F_p)^tann ⊗ K_a with
%        f_q ↦ f^a, got from a·Id_G being a monomorphism (x ↦ x^a is onto in
%        Q̄_p^*) and so an isomorphism G_{K_a} ⥲ G_a ≃ G(F_q) carrying the
%        absolute Frobenius to the relative one; the base-change functor
%        F_p → F_q read as (E, F_E) ↦ (E, F_E^a); then 10.2, G(F_q) known as
%        the sub-gerbe G_a of G(F_p) containing H.
%   83   the character-group triangle Q ← Q̄* ← Q̄* with u_a^*(λ) = v_p(λ)/a,
%        and Proposition 10.3, Manin's theorem « énoncé sous une forme plus
%        jolie et débarassé d'hypothèses restrictives parasites »: the slopes
%        of (E,F_E) over F_q are the normalised absolute values of the roots of
%        the characteristic polynomial of F_E^a divided by a, and r_i | d_i, so
%        the image in FCriso(F̄_p) is Σ E_{p_i}^{d_i'}. The margin proposes
%        instead the valuation v_q with v_q(q) = 1, dividing by nothing. The
%        typescript stops at a « Cor. » in his hand.
%
% Page 85, one leaf of another typescript, « Topos classifiant d'un
% pro-groupe »: 2.7.1 defines an operation of a strict projective system
% G = (G_i) on a set E — a family (E_i) covering E, an operation of G_i on
% E_i, and E_i the fixed points of Ker(G_j → G_i) in E_j — and gets a U-topos
% B_G by Giraud's criterion; 2.7.2 the profinite case, where an operation is a
% continuous one, i.e. one with open stabilisers; 2.7.3, that B_G is of the
% form Ĉ only when the system is essentially constant, so for a profinite
% group only when it is finite. His hand supplies « pro- », the letters G, the
% subsection numbers, and « (à gauche, disons) ».
%
% Page 86, in his hand alone: the census of the simple types — A_n and its
% intermediate isogeny classes A_{n,d}, B_n, C_n, D_n with D_n', D_n'' for n
% even ≥ 6, E_6, E_7, E_8, F_4, G_2, each with its adjoint form, and beside
% them the forms carrying a tilde — closed by « 21 cas à examiner » and two
% alternatives, of which only b) « un ... des Chevalley » is legible.
%
% Pages 88-98, « Filtration des foncteurs fibres » (his cover, page 87), his
% own hand throughout, paginated 1 to ... in his own numbering (88 = 1,
% 89 = 2, 92 = 3):
%
%   88   the definition: a filtration of a relative fibre functor F on a rigid
%        tensor category is admissible when (0) each F^{(p)}(M) is a local
%        direct factor, (i) additive, (ii) multiplicative with
%        F^{(p)}(M⊗N) = Σ_{q+r=p} F^q M ⊗ F^r N, (iii) F^p(M̌) orthogonal to
%        F^{-p}(M), (iv) locally exhaustive and separated; then gr F is a fibre
%        functor, additive, multiplicative and commuting with duality, and
%        (iii)-(iv) force F^{(0)}(1) = 1, F^{(1)}(1) = 0.
%   89   G = Aut_S(F); gr F is the fibre functor F^P of a G-torsor
%        P = Isom_S(F, gr F), the grading given by i: G_m → G^P = G'; the
%        isomorphisms respecting the filtration and inducing the identity on
%        the graded pieces form a quasi-torseur Q under a subgroup H' of G'^P;
%        what stays to be shown is that Q admits a section locally.
%   90   the group H': for a generator M, G' = Aut(F'(M)) and H' is the set of
%        g' with g'x_p ≡ x_p mod F'^{(p+1)}(M); g' decomposes along the
%        characters of G_m and h' = Σ_{n≥1} g'_n, boxed; h' nilpotent hence
%        unipotent, b' = Σ_{n≥0} g'_n parabolic, the centraliser of i a Levi —
%        and the N.B. to study one-parameter subgroups against parabolics with
%        their Levi decomposition.
%   92   the summary: for k suitable and M semi-simple, an admissible filtered
%        fibre functor is the same as (i) an admissible graded one, i.e. F'
%        together with i: G_m → G' = Aut_S F', and (ii) a right torseur Q'
%        under the uni-parabolic group of G' defined by i; two homomorphisms
%        i_1, i_2 with i_1(λ)i_2(λ) = f(λ) are called complementary, and their
%        parabolics are opposite with a common Levi in the intersection.
%   94   the Tate structure ξ: G' → G_m, with ξi_1(λ) = λ^n; the Exemple: S a
%        connected préschéma over Q, M the motives on S, the De Rham fibre
%        functor admissible, whose F' is the Hodge functor, and over a point Q
%        is trivial, so the filtered De Rham functor splits — but not
%        canonically, and not stably under the automorphisms of the functor.
%   95   S = Spec K with i: K → C: the Betti fibre functor F'', Gaga and the
%        theorem of harmonic integrals giving F''_C ≃ F_C ≃ F'_C, hence a real
%        structure with involution σ on F_C and the bigrading
%        F^{p,q}_C = F^{(p)}_C · σ(F^{(q)}_C).
%   96   the same in group terms: G affine over k, i: G_m → G, U ⊂ P ⊂ G with
%        P = Aut_{⊗,filt} T and U = Aut_{⊗,filt,id sur Gr}(T); for G reductive
%        in characteristic 0 with a faithful module V, P is the stabiliser of
%        Fil^p(V) = Σ_{i≥p} Gr^i(V) in Aut_k(V) and U the kernel of
%        P → ∏ Aut_k Gr^i(V); in general both are intersections of such;
%        g = Σ g^α along the characters of G_m, with [g^α, g^β] ⊂ g^{α+β}.
%   97   p = Σ_{α≥0} g^α and u = Σ_{α>0} g^α, computed from x V^β ⊂ Σ_{β'≥β}
%        V^{β'}; then a block, cancelled and boxed, on P = G ∩ P', U = G ∩ U'
%        and the closed immersion G/P ↪ G'/P'.
%   98   the Proposition: G connected of finite type and reductive makes P
%        parabolic and U its unipotent radical; the Corollaire that every
%        ⊗-filtration of a fibre functor over a perfect field splits, « En
%        effet, H^1(k,U) = 0 par Rosenlicht »; and a Remarque, that the uTu^{-1}
%        generate T·U, with V^{T·U} = ∩_{u∈U'} V^{uTu^{-1}} and V^G = V^P. The
%        last block is boxed and struck out by a dozen diagonals.
%
% Page 99 is a cover sheet in his hand, « Notes Saavedra », naming a run that
% begins after this batch; it takes no \page{} and is given as a heading.
%
% Pages 84, 91, 93 and 100 are blank and are absent.
%
% The handwritten run is written at speed. Its statements and formulas carry;
% the connective prose does not, and page 90 in particular is a palimpsest
% whose first six lines will not resolve. They are left \ill{} rather than
% filled. Two readings are flagged rather than repaired: page 88's insertion
% « effectif », whose point of attachment the leaf does not settle, and page
% 97's « ν + μ ≤ α i.e. μ ≤ 0 », whose sense the condition stated four lines
% above would reverse.
%
% In pages 81-83 the underlinings are the typescript's own, as in batch 4;
% in his own hand, pages 88-98, they are not used for the number fields and
% are not supplied.
%
% The dating is the inventory's own, for the whole shelfmark.
\documentclass[11pt,a4paper]{article}
% Le préambule commun à toutes les transcriptions du fonds Grothendieck.
%
% Il tient en une idée : **l'incertitude se compose, elle ne se résout pas.**
% Une transcription qui lisse un mot douteux a détruit la seule chose pour
% laquelle on la faisait — un lecteur doit pouvoir savoir, à chaque endroit,
% ce qui était sur le feuillet et ce qui a été ajouté. Les six macros qui
% suivent sont donc l'appareil critique, et non de la décoration.
%
% Les mêmes commandes sont reconnues par `scripts/render.mjs`, qui produit la
% vue de lecture du site. Une macro ajoutée ici doit l'être aussi là-bas,
% faute de quoi le rendu échouera bruyamment — ce qui est voulu : un rendu
% silencieusement faux est pire qu'un rendu absent.

\ProvidesPackage{grothendieck}[2026/08/08 Transcriptions du fonds Grothendieck]

\usepackage{amsmath,amssymb,amsthm}
\usepackage{mathrsfs}
\usepackage{stmaryrd}
% Les diagrammes commutatifs sont partout dans ce fonds : c'est la langue
% même de la géométrie algébrique de Grothendieck, et une transcription qui
% les rendrait en prose ne transcrirait rien.
\usepackage{tikz-cd}
% Français par défaut — c'est la langue des feuillets — mais l'anglais est
% chargé aussi : la traduction et le résumé sont des documents anglais, et la
% typographie française (espaces avant les deux-points) y serait une faute.
% Ces documents basculent par \selectlanguage{english} en tête de corps.
\usepackage[english,french]{babel}
\usepackage{fontspec}
\usepackage{geometry}
\geometry{a4paper,margin=25mm,marginparwidth=28mm}
\usepackage{marginnote}
\usepackage{xcolor}
% ulem, not soul. `soul`'s \st and \ul are built on a character-by-character
% scan that XeLaTeX's font machinery breaks: the document fails with an
% undefined control sequence rather than degrading. ulem does the same two jobs
% and is Unicode-engine safe. `normalem` stops it hijacking \emph, which the
% transcriptions use for Grothendieck's own emphasis.
\usepackage[normalem]{ulem}
\usepackage{hyperref}
\hypersetup{colorlinks=true,linkcolor=black,urlcolor=black,pdfborder={0 0 0}}

% --- Métadonnées du lot -----------------------------------------------------
% Reprises telles quelles par le script de rendu, qui les affiche en tête de
% la vue de lecture. Elles fixent ce que le fichier prétend couvrir : sans
% elles, une transcription flotte sans point d'attache dans un fonds de seize
% mille feuillets.

\newcommand{\folder}[1]{\gdef\grothFolder{#1}}
\newcommand{\batch}[1]{\gdef\grothBatch{#1}}
\newcommand{\leaves}[2]{\gdef\grothFirst{#1}\gdef\grothLast{#2}}
\newcommand{\foldertitle}[1]{\gdef\grothFoldertitle{#1}}
\gdef\grothFolder{?}\gdef\grothBatch{?}\gdef\grothFirst{?}\gdef\grothLast{?}\gdef\grothFoldertitle{}

% --- Le résumé --------------------------------------------------------------
%
% Il ouvre la lecture modernisée, sous le titre « Résumé », et il est composé
% en petit. Ce n'est pas un ornement : c'est le seul passage du document qui
% ne soit pas des mathématiques, et un lecteur doit pouvoir le reconnaître
% avant de l'avoir lu — puis le sauter s'il connaît déjà le sujet, ou s'y
% arrêter s'il ne le connaît pas. Le filet qui le suit marque la reprise du
% corps.

\newenvironment{resume}
  {\small\setlength{\parskip}{0.45em}}
  {\par\medskip\hrule\medskip}

% --- Mots-clés ---------------------------------------------------------------
%
% En anglais, en clôture du résumé de la lecture modernisée : le vocabulaire
% moderne sous lequel chercher ce que les feuillets construisent. Le manifeste
% du site (`npm run manifest`) les extrait pour étiqueter la cote entière —
% cette ligne est la source unique des tags, il n'y a pas de fichier à côté.
\newcommand{\keywords}[1]{\par\smallskip\noindent
  {\footnotesize\textsc{Keywords} --- \textit{#1}}\par}

% --- L'appareil critique ----------------------------------------------------

% Le numéro de feuillet, en marge. C'est l'ancre qui permet de régler un
% désaccord sur une formule en regardant *un* feuillet plutôt qu'une chemise
% de six cents. Le site s'en sert aussi pour tourner les pages du fac-similé
% quand on fait défiler la transcription.
\newcommand{\leaf}[1]{%
  \marginnote{\footnotesize\color{gray!70}\textsf{#1}}%
  \label{leaf:#1}\ignorespaces}

% Illisible : on ne devine pas. Un mot inventé qui se lit comme les autres est
% le pire résultat possible.
\newcommand{\ill}{\textcolor{red!60!black}{[\,\ldots\,]}}

% Lecture douteuse : le mot est donné, et signalé comme incertain.
\newcommand{\uncertain}[1]{\uline{#1}}

% Ajout de l'éditeur — une lettre manquante, un mot sous-entendu.
\newcommand{\add}[1]{\textcolor{blue!50!black}{[#1]}}

% Biffé par Grothendieck lui-même. Ce qu'il a rayé fait partie du document :
% une rature dit souvent où la pensée a changé de direction.
\newcommand{\struck}[1]{\sout{#1}}

% Note du transcripteur, sur l'état matériel du feuillet.
\newcommand{\note}[1]{\par\smallskip{\small\color{gray!60!black}\textit{[#1]}}\par\smallskip}

% Note marginale de Grothendieck, distincte du corps du texte.
\newcommand{\marginal}[1]{\par\smallskip\noindent\hspace{1em}%
  {\small\color{gray!60!black}#1}\par\smallskip}

% --- Titre ------------------------------------------------------------------

\AtBeginDocument{%
  \begin{center}
    {\large\bfseries Fonds Alexandre Grothendieck --- cote n\textsuperscript{o}~\grothFolder}\\[2pt]
    {\itshape\grothFoldertitle}\\[4pt]
    {\small feuillets \grothFirst--\grothLast\ (lot \grothBatch)}\\[2pt]
    {\footnotesize\color{gray!60!black}Transcription établie d'après le fac-similé numérisé
     par l'Université de Montpellier.\\
     Les lectures incertaines sont soulignées, les illisibles notées [\,\ldots\,],
     les ajouts entre crochets.}
  \end{center}
  \vspace{1em}}

\endinput


\folder{10}
\batch{5}
\pages{81}{100}
\dating{[à partir de 1958]}
\watermark{Édition de démonstration}
\foldertitle{Catégories tensorielles : notes manuscrites (s.d.), tapuscrits (s.d.)}

\begin{document}

\section*{Catégories tannakiennes définies par des cristaux, fin (pages 81 à 83)}

\note{Suite et fin du tapuscrit commencé page 66 et interrompu en bas de la
page 80 sur « D'ailleurs, ». Comme dans le lot précédent, les lettres de
groupes, les indices $a$ et plusieurs signes manquaient à la machine et sont
encrés de sa main ; on ne les relève pas un à un. Ses insertions sont données
entre crochets, ses suppressions barrées, ses notes marginales en retrait.}

\page{81}
les $E_{i}$ qui sont à composantes isopentiques $\lambda$ tels que
$0 < \lambda \leq 1$ sont exactement les modules de Dieudonné des groupes
formels $p$-divisibles sur $k$ (la catégorie de ces modules de Dieudonné
\struck{correspond} \add{est opposée} à la catégorie de ces groupes, mod
isogénie).

D'autre part, un $F$-isocristal est d'amplitude isopentique contenue dans
$[0,1]$ (intervalle fermé) sss c'est le module de Dieudonné d'un groupe de
Barsotti-Tate sur $k$ (la catégorie des $F$-isocristaux envisagée étant
\add{anti-}équivalente à la catégorie des groupes de B.T. sur $k$ mod.\ isogénie …) : le $F$-isocristal unité correspond au groupe de B.T.\ ind-étale $\underline{\mathbb{Q}}_{p}/\underline{\mathbb{Z}}_{p}$ sur $k$.

Supposant toujours $k$ parfait, la condition
\struck{\uncertain{qu'il existe}} que $(E, F_{E})$ soit effectif et qu'il
existe un $V_{i}$ \add{effectif} tel que
$F V_{i} = V_{i} F = p^{i}\,\mathrm{id}$ (cf n° 2) signifie que $E$ est
d'amplitude isopentique contenue dans $[0,i]$.
\note{« effectif » est tapé en interligne ; un trait à l'encre le rattache au
point situé entre $V_{i}$ et « tel », et le même trait traverse « tel ». La
page ne décide pas si l'insertion remplace « tel que » ou s'y ajoute.}

Il est clair alors que pour $\underline{\mathbb{F}}_{q}$ comme pour
$\underline{\mathbb{F}}_{p}$, la réponse à 6.2 est affirmative, donc les
considérations du n° 7 s'appliquent : on trouve la structure de
$\underline{G}(\underline{\mathbb{F}}_{q})$ comme gerbe $+$ ses structures
supplémentaires examinées jusqu'à présent.

Le foncteur fibre \add{canonique} \add{$(E,F_{E}) \mapsto E$} sur
$\mathrm{FCriso}(\underline{\mathbb{F}}_{q})$ \struck{à valeurs} sur le corps
$K_{a} = K(\underline{\mathbb{F}}_{q = p^{a}})$ (i.e.\ section canonique de
$\underline{G}(\underline{\mathbb{F}}_{q})$ sur $K_{a}$) est muni d'une
structure supplémentaire, savoir un automorphisme linéaire $F_{E}^{a}$, de
sorte que le groupe \struck{algébrique} sur \struck{$k$} $K_{a}$ correspondant
à cette section de $\underline{G}(\underline{\mathbb{F}}_{q})$ sur $K_{a}$,
savoir le sous-groupe $G_{a}$ de $G_{K_{a}}$ ouvert d'indice $a$
(correspondant au sous-groupe des racines $a$-èmes de $1$ dans le groupe des
caractères $\bar{\underline{\mathbb{Q}}}^{*}$ de $G_{K_{a}}$),
\marginal{\uncertain{prendre} plutôt $= f_{q}$ (comme un automorph.\ du
foncteur identique de $\mathrm{FCriso}(\underline{\mathbb{F}}_{q})$, soit donc
un élément de \uncertain{$G_{p}(\underline{\mathbb{Q}}_{p})$} …}

\page{82}
est muni d'un élément \add{$f_{q}$} de $G_{a}(K_{a})$ (« Frobenius relatif à
$\underline{\mathbb{F}}_{q}$ »). Ce dernier n'est autre d'ailleurs que
$f^{a}$, où $f$ est le générateur canonique de $G$ (le Frobenius absolu !).
\struck{on voit donc que $G_{a}$ est encore engendré par} Notons que
$a \cdot \mathrm{Id}_{G}$ est un monomorphisme (car $x \mapsto x^{a}$ est un
épimorphisme dans $\bar{\underline{\mathbb{Q}}}^{*}_{p}$), donc induit un
isomorphisme entre $G_{K_{a}}$ et
$G_{a} \simeq G(\underline{\mathbb{F}}_{q})$, transformant le frobenius
\struck{relatif} absolu en le frobenius relatif à $\underline{\mathbb{F}}_{q}$.
Grâce à ceci, on voit :

\textbf{Proposition 10.1.} La catégorie
$\mathrm{FCriso}(\underline{\mathbb{F}}_{q})^{\mathrm{tann.}}
\otimes_{\underline{\mathbb{Q}}_{p}} K_{a}$ est équivalente à la catégorie
$\mathrm{FCriso}(\underline{\mathbb{F}}_{p})^{\mathrm{tann.}}
\otimes_{\underline{\mathbb{Q}}_{p}} K_{a}$, avec $f_{q}$ correspondant à
$f^{a}$ (où on peut, si on veut, considérer $f_{q}$ et $f$ comme des
endomorphismes des foncteurs identiques de
$\mathrm{FCriso}(\underline{\mathbb{F}}_{q})$ resp.\ $\mathrm{FCriso}(\underline{\mathbb{F}}_{p})$), et correspondance des foncteurs
fibres canoniques.
\marginal{i.e.\ la catégorie des couples $(E, F'_{E})$ d'un vectoriel sur
$K_{a}$ et d'un automorphisme de ce vectoriel …)}

Le foncteur déduit du changement de base $\underline{\mathbb{F}}_{p} \to
\underline{\mathbb{F}}_{q}$
\[
\mathrm{FCriso}(\underline{\mathbb{F}}_{p})
  \otimes_{\underline{\mathbb{Q}}_{p}} K_{a}
\longrightarrow
\mathrm{FCriso}(\underline{\mathbb{F}}_{q})
  \otimes_{\underline{\mathbb{Q}}_{p}} K_{a}
\]
s'identifie, par l'équivalence précédente, au foncteur qui à tout couple
$(E, F_{E})$ d'un vectoriel sur $K_{a}$ et d'un automorphisme dudit, associe le
couple $(E, F_{E}^{a})$ ; \add{analogue \uncertain{pour}
$\mathrm{FCriso}(\underline{\mathbb{F}}_{1}) \to
\mathrm{FCriso}(\underline{\mathbb{F}}_{q})
\otimes_{\underline{\mathbb{Q}}_{p}} K_{a}$ ; il s'identifie :
$(E, F_{E}) \mapsto (E, F_{E}^{a})$.}
\note{L'indice du premier $\underline{\mathbb{F}}$ de l'ajout est un simple
trait vertical ; le contexte appellerait $p$. On laisse comme il est écrit.}

10.2. Comme on connait $\underline{G}(\underline{\mathbb{F}}_{q})$ comme une
sous-gerbe $G_{a}$ de $G = \underline{G}(\underline{\mathbb{F}}_{p})$, et que
celle-ci contient $\underline{H}$, on connait donc en principe l'inclusion
\[
\underline{H} \longrightarrow
G_{a} = \underline{G}(\underline{\mathbb{F}}_{q}) ,
\]
donc aussi le foncteur
\[
\operatorname{Rep}(G_{a}) \simeq
\mathrm{FCriso}(\underline{\mathbb{F}}_{q}) \longrightarrow
\operatorname{Rep}(\underline{H}) ,
\]
\struck{considérons} et le foncteur correspondant
\[
\mathrm{FCriso}(\underline{\mathbb{F}}_{q})
  \otimes_{\underline{\mathbb{Q}}_{p}} K_{a}
\longrightarrow
\operatorname{Rep}(\underline{H})
  \otimes_{\underline{\mathbb{Q}}_{p}} K_{a} .
\]
\note{Cette dernière phrase et la formule qui la suit sont biffées d'un long
trait à la main.}

L'effet de ces foncteurs sur les classes d'isomorphie d'éléments est connu
(comme chaque fois qu'il s'agit d'un morphisme d'une \struck{gerbe de}
catégorie tannakienne à lien commutatif dans une autre à lien de type
multiplicatif) par la connaissance de l'homomorphisme de liens correspondant

\page{83}
\[
\underline{H} \xrightarrow{\;u_{a}\;} G_{a} \longrightarrow G ,
\]
lui-même déterminé par la connaissance de l'homomorphisme en sens inverse sur
les groupes de caractères

\begin{tikzcd}
\underline{\mathbb{Q}} &
\bar{\underline{\mathbb{Q}}}^{*} \arrow[l, "u_{a}^{*}"'] &
\bar{\underline{\mathbb{Q}}}^{*} \arrow[l, "\lambda^{a}"']
  \arrow[ll, bend left=25, "|\lambda|"']
\end{tikzcd}

On lit sur ce diagramme que, pour l'isomorphisme naturel du groupe des
caractères de $G_{a}$ avec $\bar{\underline{\mathbb{Q}}}^{*}$, provenant de
l'élément $f_{q} \in G_{a}(\underline{\mathbb{Q}})$, \add{on} \struck{qui}
donne
\[
u_{a}^{*}(\lambda) = \frac{v_{p}(\lambda)}{a} .
\]

On trouve donc le résultat suivant de Manin, énoncé sous une forme plus jolie
et débarássé d'hypothèses restrictives parasites :
\note{Sic, l'accent est de la machine.}

\textbf{Proposition 10.3.} Soit $(E, F_{E})$ un $F$-isocristal sur
$\underline{\mathbb{F}}_{q}$, $q = p^{a}$, et considérons le polynôme
caractéristique $f(T)$ de $F_{E}^{a}$, soient $\lambda_{1}, \dots,
\lambda_{N}$ ($N = \operatorname{rang} E$) ses racines, et
$c_{1}, \dots, c_{N}$ leurs valeurs absolues normalisées
($c_{i} = v_{p}(\lambda_{i})$, $v_{p}(p) = 1$). Alors \struck{sur} dans
$\operatorname{Rep}(\underline{H})
\otimes_{\underline{\mathbb{Q}}_{p}} \bar{\underline{\mathbb{Q}}}_{p}
\simeq \mathrm{FCriso}(\underline{\mathbb{F}}_{q})
\otimes_{\underline{\mathbb{Q}}_{p}} \bar{\underline{\mathbb{Q}}}_{p}$,
l'image de $(E, F_{E})$ a comme pentes les nombres
$c_{1}/a, \dots, c_{N}/a$, de sorte que, si les \add{différentes} pentes
possibles sont \struck{$p_{i}$ avec multiplicités}
$p_{i} = s_{i}/r_{i}$ ($s_{i}, r_{i} \in \underline{\mathbb{Z}}$,
$r_{i} \geq 1$, $(s_{i}, r_{i}) = 1$) avec multiplicité $d_{i}$, on a
$r_{i} \mid d_{i}$ i.e.\ $d_{i} = d_{i}' r_{i}$, et l'image de $(E, F_{E})$
dans $\mathrm{FCriso}(\bar{\underline{\mathbb{F}}}_{p})$ est isomorphe à
$\sum_{i} E_{p_{i}}^{d_{i}'}$.
\marginal{N.B. ou : --- prendre plutôt la \uncertain{val.} $v_{q}$ telle que
$v_{q}(q) = 1$, i.e.\ \ill{} les valeurs propres $v_{q}(\lambda_{i})$ (sans
diviser par $a$))}

\textbf{Cor.}
\note{Le mot est écrit de sa main et souligné ; rien ne le suit, et le
tapuscrit s'arrête là.}

\section*{Topos classifiant d'un pro-groupe (page 85)}

\note{Feuillet d'un tout autre tapuscrit, corrigé de sa main : c'est lui qui
ajoute « pro- », les lettres de groupes, les numéros de sous-section et
« (à gauche, disons) ».}

\page{85}
\textbf{2.7. Topos classifiant d'un \add{pro-}\struck{profini} groupe.}

2.7.1 Soit \add{$\underline{G} =$} $(G_{i})_{i \in I}$ un système projectif
\struck{strict} de groupes, avec $G_{i}, I \in \mathcal{U}$. On suppose le
système projectif strict, i.e.\ les morphismes de transition
$G_{j} \to G_{i}$ surjectifs. Si $E$ est un ensemble, on
\struck{dit que $\underline{G}$ opère sur $E$} appelle opération de
$\underline{G}$ sur $E$ \add{(à gauche, disons)} la structure suivante :
a) une famille $(E_{i})_{i \in I}$ de parties de $E$, de réunion $E$ ;
b) pour tout $i \in I$, une opération du groupe $G_{i}$ sur l'ensemble
$E_{i}$ ; on suppose de plus ces données soumises à la condition suivante :
pour $j \geq i$, $E_{i}$ est le sous-ensemble de $E_{j}$ formé des éléments
fixes sous le groupe noyau de $G_{j} \to G_{i}$. On dit \add{aussi} que
$\underline{G}$ opère sur $E$ (à gauche) si on s'est donnée une opération de
$\underline{G}$ sur $E$ (à gauche). Les ensembles $\in \mathcal{U}$ munis
d'une opération de $\underline{G}$ forment une catégorie de façon évidente.
On constate aussitôt, grâce au critère de Giraud, que cette catégorie est un
\add{$\mathcal{U}$-}topos. On le note $B_{\underline{G}}$ et on l'appelle
topos classifiant de $\underline{G}$. Lorsque $I$ admet un objet
\struck{\ill{}} initial $i_{0}$, posant $G = G_{i_{0}}$, on retrouve le topos
classifiant de 2.3.

\add{2.7.2} Un autre exemple important est celui où les groupes $G_{i}$ sont
finis, de sorte que
\[
G = \varprojlim G_{i}
\]
est un groupe topologique compact totalement discontinu, ou groupe profini.
Une opération de $\underline{G}$ sur $E$ revient alors à une opération de $G$
sur $E$ qui est continue, ou ce qui revient au même, telle que le
stabilisateur de tout point de $E$ est un sous-groupe ouvert de $G$. Le topos
classifiant $B_{\underline{G}}$ sera aussi noté $B_{G}$, où, bien entendu,
$G$ doit être considéré comme muni de sa topologie profinie.
\marginal{à la ligne}

\add{2.7.3.} Il est facile de vérifier, \struck{en} utilisant les remarques de
2.6, que le topos $B_{\underline{G}}$ défini par un système projectif strict
\struck{de} $\underline{G} = (G_{i})_{i \in I}$ de groupes n'est équivalent à
un topos de la forme $\hat{C}$ que si ce système projectif est essentiellement
constant ; dans le cas d'un groupe profini, cela signifie que ce groupe est en
fait fini.

\section*{Les types simples, un recensement (page 86)}

\note{De sa main seule, au dos du feuillet précédent. Les formes de la colonne
de droite portent un tilde ; l'exposant « ad » est celui de la forme adjointe.}

\page{86}
\begin{itemize}
  \item $A_{n}$, $A_{n}^{\mathrm{ad}}$, $(A_{n,d})$ ($d \mid n$,
        $1 \leq d < n$) ; en regard : $\tilde{A}_{n}$,
        $\tilde{A}_{n}^{\mathrm{ad}}$, $(\tilde{A}_{n,d})$
  \item $B_{n}$, $B_{n}^{\mathrm{ad}}$
  \item $C_{n}$, $C_{n}^{\mathrm{ad}}$
  \item $D_{n}$, $D_{n}^{\mathrm{ad}}$, $D_{n}'$, $D_{n}''$ (si $n$ pair
        $\geq 6$) ; en regard : $\tilde{D}_{n}$,
        $\tilde{D}_{n}^{\mathrm{ad}}$, \uncertain{$\tilde{D}_{n}'$}
  \item $E_{6}$, $E_{6}^{\mathrm{ad}}$ ; en regard : $\tilde{E}_{6}$,
        $\tilde{E}_{6}^{\mathrm{ad}}$
  \item $E_{7}$, $E_{7}^{\mathrm{ad}}$
  \item $E_{8}$
  \item $F_{4}$.
  \item $G_{2}$
\end{itemize}

21 cas à examiner
\begin{itemize}
  \item a) un produit \ill{} des groupes \ill{}
  \item b) un \ill{} des Chevalley
\end{itemize}

\section*{Filtration des foncteurs fibres (pages 88 à 98)}

\note{Le titre est de sa main, sur un feuillet par ailleurs vierge qui précède
la suite (page 87). Tout ce qui suit est de sa main. Il numérote lui-même les
feuillets : la page 88 porte son n° 1, la page 89 son n° 2, la page 92 son
n° 3. Il n'y souligne pas les corps de nombres et on ne le supplée pas.}

\page{88}
\emph{Filtration sur foncteur fibre pour catégorie tensorielle}

Soit \struck{$E$} $\mathcal{M}$ une catégorie tensorielle \struck{sur}
« rigide \ill{} » sur un corps $k$, $S$ un schéma sur $k$, $F$ un foncteur
fibre relatif :
\[
\mathcal{M} \longrightarrow
\text{faisceaux localement libres sur } S .
\]

Une filtration de $F$ par des sous-foncteurs $F^{(p)}$ est dite
\emph{admissible} si elle satisfait aux conditions :
\begin{itemize}
  \item (0) Chaque $F^{(p)}(M)$ est loc.\ facteur direct dans $M$
        ($\forall M \in \mathcal{O}b\,\mathcal{M}$)
  \item (i) Additivité : $F^{(p)}(M \ast N) = F^{(p)} M \ast F^{(p)} N$
  \item (ii) Multiplicativité :
        $F^{(p)}(M \otimes N) = \sum_{q + r = p} F^{q}(M) \otimes F^{r}(N)$
  \item (iii) $F^{p}(\check{M})$ est l'orthogonal de \uncertain{$F^{-p}(M)$}
  \item (iv) Pour chaque $M$, la filtration $F^{(p)}(M)$ est loc.\ exhaustive
        et séparée.
\end{itemize}
\note{Le signe de (i) est une petite étoile à six branches, donnée ici
$\ast$ ; l'exposant de (iii) est repassé sur la page.}

Considérons donc le foncteur
\[
M \rightsquigarrow \operatorname{gr} F(M) .
\]
En vertu de (0), il est à valeurs dans les modules loc.\ cohérents
\add{(\uncertain{p.\ libres})} sur $S$, (i) indique qu'il est additif, (ii)
qu'il est multiplicatif, et (iii) qu'il commute à : $\vee$.
\note{Le signe final est le petit « v » de la dualité, celui que (iii) écrit
$\check{M}$.}

N.B. \struck{Ses conditions impliquent} (iii) et (iv) \uncertain{exigent}
$F^{(0)}(\mathbf{1}) = \mathbf{1}$, $F^{(1)}(\mathbf{1}) = 0$, d'où
$\operatorname{gr} F(\mathbf{1}) = \mathbf{1}$ \ill{} en degré $0$.
Si \struck{\ill{}} $\mathcal{M}$ ½simple, alors $\operatorname{gr} F$
\emph{exact} (sinon il faut en mettre une condition dans (i).)

\page{89}
\note{Son n° 2.}
\[
G = \underline{\operatorname{Aut}}_{S}(F)
\]
(schéma en groupes affine sur $S$), alors $\operatorname{gr}(F)$, en tant que
foncteur-fibre (à un \uncertain{oubli} \ill{}), est défini par un
$G$-torseur $P$, \add{$\operatorname{gr}(F) \simeq F^{P}$,
$P = \underline{\operatorname{Isom}}_{S}(F, \operatorname{gr} F)$.}
(De plus, la graduation sur \struck{$\otimes$} le foncteur \struck{sur}
$\operatorname{gr}(F)$ est définie par un homomorphisme
\[
\mathbb{G}_{m_{S}} \xrightarrow{\;i\;} G^{P} = G' . )
\]

Toutes ces constructions sont \uncertain{locales} et compatibles avec
changement de base. Parmi les isom.\ de $F$ sur $\operatorname{gr}(F)$, on
distingue ceux \add{et induisant l'identité sur les gradués associés,} qui
respectent la filtration. Ils forment (\ill{}) un sous-foncteur $Q$ de $P$,
qui est un \add{quasi-}torseur à g.\ sous le sous-groupe $H'$ de $G'^{P}$
\add{et induisant l'identité sur $\operatorname{gr}$,} qui respectent la
filtration \add{canonique} de $F' = \operatorname{gr} F$, et une torseur à
droite sous le sous-groupe analogue $H$ de $\underline{Q}$.
\add{bien \uncertain{puisque} c'est un torseur,}

Cependant, il faut démontrer que loc.\ $Q$ admet une section, donc qu'il
existe un isom.\ respectant les filtrations de $F$ et $F'$ et induisant
l'identité sur les gr.\ associés.
\marginal{(en supposant $\mathcal{M}$ de type fini)}

\page{90}
On peut \ill{} que $F$ et $F'$ \uncertain{isomorphes}, \struck{\ill{}} s'il est
vrai que deux foncteurs fibres sont isomorphes \uncertain{loc.} \ill{}, du
moins \struck{en tout cas} si le \ill{} de \ill{}, \add{en supposant que
$\mathcal{M}$ ½simple} (\ill{} comme lien de \uncertain{schémas} \ill{} et plat
sur $S$, \ill{}).
\note{Les six premières lignes de la page sont écrites à toute vitesse, sur un
passage déjà repris ; seuls les mots donnés se lisent.}

Déterminons le groupe $H'$, et pour cela choisissons un \emph{générateur} $M$,
alors $G' = \underline{\operatorname{Aut}}_{\mathcal{O}_{S}}(F'(M))$, et $H'$
est le sous-groupe de $G'$ formé des $g'$ tels que,
$\forall p \in \mathbb{Z}$,
\[
g' x_{p}^{(p)} \equiv x_{p}^{(p)} \qquad
\bigl( F'^{(p+1)}(M) \bigr)
\]
\note{Le membre de droite est repris et biffé ; l'expression entre parenthèses
est portée au-dessus du passage biffé, et un « $\in$ » clôt la ligne sans
complément.}
\marginal{Construction intrinsèque de $G'^{\mathrm{ad}}$,
$\mathbb{G}_{m} \xrightarrow{\;i\;} G'$ (\uncertain{comme pour} $G$).}

Les transformations infinitésimales de $H'$ s'obtiennent ainsi : comme \ill{}
de $\mathfrak{g}$ \ill{}, i.e.\ $i : \mathbb{G}_{m} \to G'$,
$\mathfrak{g}'$ se décompose suivant les caractères de $\mathbb{G}_{m}$, et on
a
\[
\boxed{\;\mathfrak{h}' = \sum_{n \geq 1} \mathfrak{g}'_{n}\;}
\]

N.B. Le sous-groupe qui \uncertain{respecte} la filtration, \ill{} :
l'algèbre de Lie $\mathfrak{b}' = \sum_{n \geq 0} \mathfrak{g}'_{n}$. C'est le
produit semi-direct \ill{} de $\mathfrak{g}$, et de $\mathfrak{h}'$.
$\mathfrak{h}'$ est nilpotente, et correspond d'ailleurs à un groupe
\emph{unipotent}, tandis que $\mathfrak{b}'$ est parabolique ; le
centralisateur de $i$ est un sous-groupe de Levi.
\note{« groupe » est barré et « unipotent » repris au-dessus.}

N.B. Étudier les relations entre « ss-gpes à $1$ paramètre » et sous-groupes
paraboliques avec déc.\ de Levi.
\marginal{N.B. $H'$ \ill{} \emph{gpes connexes} (étant \emph{unipotent} !)}

\page{92}
\note{Son n° 3.}
Bien entendu, on \uncertain{récupère} $F$, \uncertain{donc aussi} la
filtration de $F$ connaissant $F' = F^{P}$,
$i : \mathbb{G}_{m} \to G' = G^{P}$
$\bigl( = \underline{\operatorname{Ad}}_{S}(P) \bigr)$, et le
sous-$H'$-torseur \add{\uncertain{à gauche}} $Q$ de $P$.

En \emph{résumé} (si $k$ est \ill{}, $\mathcal{M}$ ½simple), la donnée d'un
\struck{foncteur} \add{fibre} \emph{filtré admissible} sur $\mathcal{M}$, à
« valeurs dans $S$ », équivaut à la donnée
\begin{itemize}
  \item (i) d'un foncteur fibre \emph{gradué admissible} sur $\mathcal{M}$,
        \add{$F'$}, à valeurs dans $S$, i.e.\ d'un foncteur fibre $F'$ sur
        $\mathcal{M}$, et d'un hom.\ $i : \mathbb{G}_{m_{S}} \longrightarrow G' =
        \underline{\operatorname{Aut}}_{S} F'$ ;
  \item (ii) d'un \struck{$G'^{\geq}$}-torseur \struck{$P$} $Q'$ à droite,
        \ill{} \struck{$H'$} \struck{groupe} uni-parabolique de $G'$ défini par
        $i$.
\end{itemize}

À ces données, on associe le foncteur fibre $F = F'^{Q'}$, avec la filtration
déduite par descente de la filtration \struck{canonique} de $F'$
\struck{gradué} \ill{} sa graduation. D'ailleurs $Q$ \ill{} la filtration du
foncteur $F$ « splitte » …

Si $\mathcal{M}$ est muni \struck{des} \add{de la} donnée \ill{} et des
structures \uncertain{intrinsèques}, \ill{} dans \ill{} :
$i : \mathbb{G}_{m_{S}} \to G'$ \ill{} \uncertain{central}, et : $i_{1}$, il
convient de prendre $i_{2}$ défini par
\[
i_{1}(\lambda)\, i_{2}(\lambda) = f(\lambda)
\]
Alors $i_{1}$, $i_{2}$ sont dits \emph{complémentaires}. Les paraboliques
associés sont \emph{opposés} (\struck{\ill{}}), et leur intersection est un
Levi commun.

\page{94}
s'il y a aussi la structure intrinsèque de Tate, définissant donc
\[
\xi : G' \longrightarrow \mathbb{G}_{m}
\]
alors $\xi\, i_{1}(\lambda) = \lambda^{n}$, et l'entier $n$ est le
\uncertain{degré} du \struck{\ill{}} « \uncertain{motif} de Tate »
\uncertain{effectif} (\uncertain{noté} $T$ \uncertain{ou}
$(\mathbb{G}_{m})^{-1}$).

Dans le cas auquel nous nous intéressons, i.e.\ dans \ill{}, \ill{} le cas est
qu'il y a compatibilité avec la structure de Tate. --- \uncertain{Donc} :
\[
\xi\, i_{1}(\lambda) = \lambda
\]

\emph{Exemple}. Soit $S$ préschéma \add{connexe} sur $\underline{\mathbb{Q}}$,
et $\mathcal{M}$ la catégorie des motifs sur $S$. Le \emph{foncteur de De
Rham} sur $\mathcal{M}$ est donc à valeurs dans les \struck{Modules} loc.\ libres sur $S$, et est \uncertain{admissible}.

À ce titre, il définit donc une donnée $F'$ (qui n'est autre que le
\emph{foncteur de Hodge}) \struck{\ill{}} $G'$,
$i : \mathbb{G}_{m} \longrightarrow G'$, (d'où $\mathfrak{g}'^{\geq}$) et un
$G'^{\geq}$-torseur $Q$. (Étant donné la structure de $G'^{\geq}$
(\uncertain{par univ.}), il s'ensuit d'ailleurs que si \emph{$S$ est réduit à
un pt}, $Q$ est \emph{trivial}, donc le \emph{foncteur de De Rham filtré
splitte} (mais le \emph{splitting} du foncteur n'est pas canonique, en
particulier n'est pas stable par les autom.\ de ce foncteur \add{ou par ceux
qui \uncertain{respectent} la filtration},

\page{95}
et induisant l'identité sur les gradués associés.

Supposons maintenant $S = \operatorname{Spec} K$, $K$ un corps, et
considérons un hom.\ $i : K \to \mathbb{C}$. Alors la donnée de $i$ définit le
foncteur-fibre \emph{de Betti}, soit $F''$, à valeurs dans $\mathbb{Q}$.

De plus, on a \struck{des} deux isom.\ canoniques
\[
F''_{\mathbb{C}} \;\overset{\text{Gaga}}{\simeq}\; F_{\mathbb{C}}
\;\overset{\substack{\text{th. de Hodge}\\ \text{des intégrales harmoniques}}}
{\simeq}\; F'_{\mathbb{C}}
\]

D'ailleurs, l'isom.\ $F_{\mathbb{C}} \simeq F'_{\mathbb{C}} \simeq
\operatorname{gr}(F_{\mathbb{C}})$ est déduit de l'isom.\ $F_{\mathbb{C}} \simeq F''_{\mathbb{C}}$, \uncertain{traduisant} une structure
\emph{réelle} sur $F_{\mathbb{C}}$, donnée par une involution $\sigma$, et la
\emph{filtration} \add{conj.} de $F_{\mathbb{C}}$, par
\[
F^{p,q}_{\mathbb{C}} = F^{(p)}_{\mathbb{C}} \cdot
\sigma\bigl( F^{(q)}_{\mathbb{C}} \bigr) \dots
\]

\page{96}
Soit $G$ un schéma en gr.\ affine sur un corps $k$,
$i : \mathbb{G}_{m_{k}} \longrightarrow G$ un hom., i.e.\ une structure de
$\otimes$-graduation sur le foncteur \uncertain{oubli} sur
$\operatorname{Rep}(G)$.

Cette $\otimes$-graduation définit une $\otimes$-filtration, et on peut
regarder les \uncertain{$k$-groupes}
\[
U \subset P \subset G \qquad
P = \underline{\operatorname{Aut}}_{\otimes,\,\mathrm{filt}}\, T , \quad
U = \underline{\operatorname{Aut}}_{\otimes,\,\mathrm{filt},\,
\mathrm{id.\ sur\ Gr}}(T)
\]

Lorsque $G$ \add{est réductif sur $k$ de car.\ $0$, et que} admet un module
fidèle $V$, \struck{\ill{}} alors $V$ est gradué par
$\operatorname{Gr}^{i}(V)$, filtré par
\[
\operatorname{Fil}^{p}(V) = \sum_{i \geq p} \operatorname{Gr}^{i}(V)
\]
\note{Sous le $\operatorname{Gr}$ subsiste un indice biffé que la page ne
détermine pas.}
et $P$ \struck{\ill{}} est le $k$-groupe de
$\underline{\operatorname{Aut}}_{k}(V)$ laissant invariante \add{cette}
filtration, et $U$ le sous-gpe noyau de
$P \to \prod_{i} \underline{\operatorname{Aut}}_{k}
\operatorname{Gr}^{i}(V)$.

Dans le cas général, $P$ et $U$ sont des intersections de tels gr.\ pour
$P_{V}$, $U_{V}$ ($V \in \mathcal{O}b \operatorname{Rep}_{k}(G)$) ; si $P$
\ill{}, il y a un $V$ tel que $P_{V} = P$, $U_{V} = U$.

\struck{Précisons} dans $\mathfrak{g} = \operatorname{Lie} G$, $\xi \in
\mathfrak{g}$ définit l'image du générateur \struck{can.} de
$\mathfrak{t} = \operatorname{Lie} \mathbb{G}_{m}$ par
$\operatorname{Lie} i : \mathfrak{t} \to \mathfrak{g}$, et
$\mathfrak{g} = \sum_{\alpha \in \mathbb{Z}} \mathfrak{g}^{\alpha}$ la
décomposition de $\mathfrak{g}$ suivant les \uncertain{car.} de
$\mathbb{G}_{m}$ opérant sur $\mathfrak{g}$,
$V = \sum_{\beta \in \mathbb{Z}} V^{\beta}$ de \uncertain{même}. On a donc
\[
[\mathfrak{g}^{\alpha}, \mathfrak{g}^{\beta}] \subset
\mathfrak{g}^{\alpha + \beta}
\qquad
\mathfrak{g}^{\alpha} \cdot V^{\beta} \subset V^{\alpha + \beta}
\]

L'algèbre de Lie $\mathfrak{p}$ de $P_{V}$ est formée des

\page{97}
$x \in \mathfrak{g}$ tels que
\[
x V^{\beta} \subset \sum_{\beta' \geq \beta} V^{\beta'}
\qquad \text{pour tout } \beta \in \mathbb{Z} .
\]
Il suffit que l'on ait
\[
x \in \sum_{\alpha \geq 0} \mathfrak{g}^{\alpha} ,
\]
et c'est aussi nécessaire si $V$ est un module fidèle, on \uncertain{écrit}
$x = \sum_{\alpha \in \mathbb{Z}} x^{\alpha}$, et si \struck{$\alpha$} $\mu$
est le plus petit des $\alpha$ tels que $x^{\alpha} \neq 0$, i.e.\ $x = x^{\mu} + \sum_{\lambda > \mu} x^{\lambda}$, $x^{\mu} \neq 0$ :
\struck{\ill{}} $\exists\, \alpha$ tel.\ que
$x^{\mu} \cdot V^{\alpha} \neq \emptyset$, donc et comme
\[
x^{\mu} V^{\alpha} \subset V^{\alpha + \mu}
\]
il faut qu'on ait $\nu + \mu \leq \alpha$ i.e.\ $\mu \leq 0$.
\note{Les deux inégalités sont écrites « $\leq$ » sur la page, et le premier
terme se lit $\nu$ ; la condition d'invariance donnée quatre lignes plus haut
demanderait $\alpha + \mu \geq \alpha$, donc $\mu \geq 0$, en accord avec
$x \in \sum_{\alpha \geq 0} \mathfrak{g}^{\alpha}$. On laisse comme il est
écrit.}
On voit \ill{} de même que si $V$ est fidèle, on a
$\mathfrak{u} = \operatorname{Lie} U = \sum_{\alpha > 0}
\mathfrak{g}^{\alpha}$. Par passage à la limite, on \uncertain{obtient}
aisément, dans tous les cas
\[
\mathfrak{p} = \sum_{\alpha \geq 0} \mathfrak{g}^{\alpha} , \qquad
\mathfrak{u} = \sum_{\alpha > 0} \mathfrak{g}^{\alpha} .
\]

\struck{Supposons \ill{} comme \ill{} que $P = P_{V}$, $U = U_{V}$, \ill{}
donc, si $V$ \ill{} fidèle \ill{} $P' \subset P$ \ill{} réciproque \ill{}
$P' = \underline{\operatorname{Aut}}(V)$, et $U'$ le \ill{} qui stabilise la
filtration de $V$, et $U'$ le \ill{} qui induit id sur
$\operatorname{Gr}(V)$,}
\[
P = G \cap P' , \qquad U = G \cap U'
\]
\struck{par suite on a une imm.\ fermée}
\[
G/P \hookrightarrow G'/P'
\]
\note{Tout le bas de la page est enfermé dans une grande boucle tracée à la
main et traversé de longs traits ; les deux formules relevées s'y lisent
seules avec certitude.}

\page{98}
\textbf{Proposition.} Supposons \add{$\underline{G}$ connexe, de t.f., et}
\struck{$k$ parfait \ill{}} $G$ réductif. Alors $P$ est un
\uncertain{sous-gr.} \uncertain{de type} \emph{parabolique},
\struck{\ill{}} $U$ son \emph{radical unipotent},
\add{et leurs formations \uncertain{commutent} : \struck{$U$ est \ill{} c.-à-d.
à b.n.}} \struck{Dans tous les cas, la \ill{} et \ill{}}
$\Longrightarrow$ $U$ unipotent …
\marginal{Vérifier que $G$ \uncertain{sans} \uncertain{connexité} \ill{} ?}

Supposons d'abord $G$ réductif. En car.\ $0$, l'assertion résulte aisément de
la théorie \ill{}. En car.\ $p > 0$, \struck{on \ill{} que \ill{} des voisins
de \ill{} et que \ill{}} il faut sans doute travailler un peu …

\struck{Formation de} \textbf{Corollaire.} Toute $\otimes$-filtration d'un
foncteur fibre sur un corps parfait \add{[\uncertain{ou même} sur un
\struck{\ill{}} quelc.]} splitte …

\struck{\emph{Démonstration}.} En effet, $H^{1}(k, U) = 0$ par
Rosenlicht …

\emph{Remarque}. Soit $T = \operatorname{Im} i$, et soit $U'$ une \ill{} dans
$U$. Alors \struck{(si $k$ \ill{})} l'ensemble des $u T u^{-1}$
($u \in U'$) engendre \struck{\ill{}} le groupe \struck{\ill{}} $T \cdot U$
(produit \uncertain{semi-direct}) \add{[En \ill{}, \ill{} l'algèbre de Lie
$\underline{\mathfrak{a}}$ contient les \ill{} $u \in U$, donc [en car.\ $0$]
donc les \ill{} $[\mathfrak{u}, x]$, \ill{}, donc $\mathfrak{u}$]} les
$\mathfrak{g}^{\alpha}$, $\alpha > 0$ [car
$[\mathfrak{g}^{\alpha}, x] = \alpha\, \mathfrak{g}^{\alpha} =
\mathfrak{g}^{\alpha}$], donc $\mathfrak{u}$.

D'où \ill{} pour toute représentation $V$ de $G$,
\[
V^{T \cdot U} \left[ \; = \bigcap_{u \in U'} V^{u T u^{-1}} \right]
\]
\[
V^{G} = V^{P} , \qquad P \text{ étant parabolique,}
\]

\struck{\ill{} $V^{P} = V^{T \cdot U}$ \ill{} $V^{U}$ \ill{} $V$ sur $P/U$
\ill{} $V^{P/U} = V^{T}$ \ill{} $\operatorname{Lie} G$ \ill{} rang $1$ \ill{}
d'ailleurs \ill{}}
\note{Ce dernier bloc est enfermé dans un cadre et barré d'une dizaine de
longues diagonales ; seules les formules relevées se lisent, la prose qui les
relie est perdue.}

\section*{Notes Saavedra (page 99)}

\note{Ces deux mots, de sa main sur un feuillet par ailleurs vierge, nomment
la suite : le tapuscrit qu'ils annoncent commence après ce lot.}

\end{document}
